\chapter*{Remerciements}
\addcontentsline{toc}{chapter}{Remerciements}

% Police Times New Roman
\fontfamily{ptm}\selectfont
\linespread{1.0}

% Ligne horizontale
\noindent\rule{\textwidth}{2pt}

% Texte principal
\vspace{0.3cm}

{\large Je tiens à exprimer ma profonde gratitude.}

\vspace{0.3cm}

Un grand merci à mon professeur encadrant, \textbf{M. Alhoussem Chtioui}, pour le temps qu’il a consacré avec tant de soin à chaque étape de ce projet. Grâce à lui, les moments difficiles sont devenus plus clairs, car il savait toujours poser les bonnes questions au moment juste. Ce n’est pas seulement son savoir qui a compté, mais surtout sa manière calme et précise de guider sans imposer. Chaque rencontre ouvrait une nouvelle direction, jamais en ligne droite, souvent par des chemins détournés mais utiles. Il ne donnait pas les réponses toutes faites, plutôt des indices discrets qui poussaient à chercher encore. Sa façon de travailler m’a appris à mieux organiser mes idées, à vérifier deux fois avant d’affirmer. Loin des discours vides, il insistait sur l’exactitude, même dans les détails invisibles. Sans cette présence attentive, ce mémoire aurait pris une autre forme, moins solide, certainement moins aboutie. Ce que j’ai gagné va bien au-delà du sujet traité - c’est une manière différente de penser, plus exigeante, plus lucide.

\vspace{0.3cm}

Un grand merci à mon tuteur en entreprise, \textbf{M. Amine Ben Slimene}, pour son aide constante, sa grande écoute et ses indications claires durant tout le stage. Ce n’est pas tous les jours qu’on croise quelqu’un qui allie savoir-faire technique et vue d’ensemble comme il l’a fait. Grâce à lui, j’ai vu comment fonctionnent vraiment les choses derrière un projet CRM. Les idées abstraites vues en cours prennent soudain sens quand on les met face aux situations du terrain. Chaque discussion avec lui ouvrait une nouvelle porte sur ce que signifie travailler dans ce secteur. L’écart entre la salle de classe et l’usine se rétrécit vite quand on est guidé ainsi. Sa manière calme mais précise de répondre m’a évité bien des détours inutiles. Jamais je n’aurais imaginé voir autant de détails entrer en jeu dans chaque décision. Il ne donnait pas simplement des réponses, il montrait où chercher. Tout cela change profondément la façon dont je vois maintenant certaines étapes du travail. Sans cette présence attentive, beaucoup serait resté flou. La théorie garde toujours un goût partiel tant qu’elle ne touche pas la pratique. C’est exactement ce qu’il a permis : rapprocher deux mondes souvent séparés.

\vspace{0.3cm}

Un grand merci à \textbf{M. Firas Beguiga}, toujours présent, qui a guidé chaque étape avec attention. Grâce à lui, les idées ont pris forme, pas à pas. Ses remarques précises ont souvent ouvert de nouvelles pistes. Quand les difficultés techniques surgissaient, il prenait le temps d’en parler, calmement. Chaque échange apportait une pièce manquante au puzzle. Ce sont ses propositions simples mais justes qui ont fait avancer les choses. Même quand la pression montait, sa clarté d’esprit restait intacte. Loin des grandes phrases, ses interventions ciblées ont marqué toute la durée du projet. Sans ces moments d’échanges francs, bien des améliorations n’auraient pas eu lieu. Son regard acéré sur les zones fragiles a évité bien des détours. À plusieurs reprises, c’est par son calme que la direction s’est redressée. Ce travail en porte forcément la trace.

\vspace{0.3cm}

Merci beaucoup aux enseignants de l’\textbf{ISSAT de Sousse} - leur façon nette d’expliquer tient une grande place dedans. Souvent présents quand on avait besoin, ils ont donné les outils juste nécessaires pendant toute la formation d'ingénieur. Leur rigueur au quotidien comptait énormément. Chaque détail vu en cours s’est révélé utile plus tard. Ce qu’ils transmettent va bien au-delà des programmes. Sans vraiment y penser sur le moment, chacune de leurs attitudes a aidé à construire un socle solide, prêt à servir maintenant.

\vspace{0.3cm}
Merci beaucoup à l’équipe technique de \textbf{VisioAd}, qui a tout de suite ouvert la porte. Dès le départ, leur façon de s’entraider simplifiait tout. Quand un problème surgissait, ils étaient là sans attendre. Grâce à eux, avancer était naturel. Chaque échange donnait une idée utile, pas juste des mots en trop. Un frisson court le long de l'instant, reste accroché. Pas besoin de chercher ailleurs pour sentir ce poids léger entre les mains.

\vspace{0.3cm}

Voilà ceux qui rendent les choses vraies, en lisant mot après mot, sans courir. Pas besoin de dire merci : leurs yeux restent là, tranquilles, fixés. Sans ça, tout s’effondrerait, lentement, comme du papier mouillé. Un souffle retenu entre deux phrases dit autant qu’un cri parfois. Ce moment volé au reste pèse lourd, même si personne ne le mesure.

\vspace{0.5cm}

\hfill{\itshape AmenAllah}

% Retour à la configuration normale
\linespread{1.0}